% GENERATED FILE. Edit canonical lessons, then run npm run generate:books.
% canonical-source-hash: fnv1a64:0daf588476972b9b
% canonical-lessons: GE-C218-lecker, GE-C218-lieb, GE-C218-locker, GE-C218-mager, GE-C218-mutig

\chapter{Tasty, Kind, Relaxed, Thin, Brave}
\label{ch:ge-tasty-kind-relaxed-thin-brave}
\hlchaptermodality{218}

\begin{chapteropening}
\textbf{By the end of this chapter:} \emph{I can say tasty, kind, relaxed, thin and brave.}

Complete the last lesson of chapter 218: i can say tasty, kind, relaxed, thin and brave.
\end{chapteropening}

\section[lecker]{lecker — tasty}

\label{lesson:GE-C218-lecker}



\begin{quote}
Before the new one: say the German for strong, then the German for single.
\end{quote}

\subsection*{You'll want to know: lecker}
\textbf{lecker} — \textquotedblleft{}tasty\textquotedblright{}.

\begin{grammarlens}[title={What you've built}]
One more word for this chapter.
\end{grammarlens}

\subsection*{Your turn}
\begin{itemize}
\raggedright
  \item \emph{Say it:} \emph{lecker}
  \item \emph{Say it:} \emph{lecker}, once more
  \item \emph{Say it:} say \emph{lecker}
  \item \emph{Recall:} say the German for strong, then the German for single, then say \emph{lecker} again
\end{itemize}

\begin{tcolorbox}[breakable,colback=teal!4,colframe=teal!35!black,title={Before you move on}]
What does lecker mean? (\textquotedblleft{}Tasty\textquotedblright{}.) Say it once more.
\end{tcolorbox}
\section[lieb]{lieb — kind, dear}

\label{lesson:GE-C218-lieb}



\begin{quote}
Before the new one: say the German for single, then the German for tasty.
\end{quote}

\subsection*{You'll want to know: lieb}
\textbf{lieb} — \textquotedblleft{}kind, dear\textquotedblright{}.

\begin{grammarlens}[title={What you've built}]
One more word for this chapter.
\end{grammarlens}

\subsection*{Your turn}
\begin{itemize}
\raggedright
  \item \emph{Say it:} \emph{lieb}
  \item \emph{Say it:} \emph{lieb}, once more
  \item \emph{Say it:} say \emph{lieb}
  \item \emph{Recall:} say the German for single, then the German for tasty, then say \emph{lieb} again
\end{itemize}

\begin{tcolorbox}[breakable,colback=teal!4,colframe=teal!35!black,title={Before you move on}]
What does lieb mean? (\textquotedblleft{}Kind, dear\textquotedblright{}.) Say it once more.
\end{tcolorbox}
\section[locker]{locker — relaxed, loose}

\label{lesson:GE-C218-locker}



\begin{quote}
Before the new one: say the German for tasty, then the German for kind.
\end{quote}

\subsection*{You'll want to know: locker}
\textbf{locker} — \textquotedblleft{}relaxed, loose\textquotedblright{}.

\begin{grammarlens}[title={What you've built}]
One more word for this chapter.
\end{grammarlens}

\subsection*{Your turn}
\begin{itemize}
\raggedright
  \item \emph{Say it:} \emph{locker}
  \item \emph{Say it:} \emph{locker}, once more
  \item \emph{Say it:} say \emph{locker}
  \item \emph{Recall:} say the German for tasty, then the German for kind, then say \emph{locker} again
\end{itemize}

\begin{tcolorbox}[breakable,colback=teal!4,colframe=teal!35!black,title={Before you move on}]
What does locker mean? (\textquotedblleft{}Relaxed, loose\textquotedblright{}.) Say it once more.
\end{tcolorbox}
\section[mager]{mager — thin, lean}

\label{lesson:GE-C218-mager}



\begin{quote}
Before the new one: say the German for kind, then the German for relaxed.
\end{quote}

\subsection*{You'll want to know: mager}
\textbf{mager} — \textquotedblleft{}thin, lean\textquotedblright{}.

\begin{grammarlens}[title={What you've built}]
One more word for this chapter.
\end{grammarlens}

\subsection*{Your turn}
\begin{itemize}
\raggedright
  \item \emph{Say it:} \emph{mager}
  \item \emph{Say it:} \emph{mager}, once more
  \item \emph{Say it:} say \emph{mager}
  \item \emph{Recall:} say the German for kind, then the German for relaxed, then say \emph{mager} again
\end{itemize}

\begin{tcolorbox}[breakable,colback=teal!4,colframe=teal!35!black,title={Before you move on}]
What does mager mean? (\textquotedblleft{}Thin, lean\textquotedblright{}.) Say it once more.
\end{tcolorbox}
\section[mutig]{mutig — brave}

\label{lesson:GE-C218-mutig}



\begin{quote}
Before the new one: say the German for relaxed, then the German for thin.
\end{quote}

\subsection*{You'll want to know: mutig}
\textbf{mutig} — \textquotedblleft{}brave\textquotedblright{}.

\begin{grammarlens}[title={What you've built}]
One more word for this chapter.
\end{grammarlens}

\subsection*{Your turn}
\begin{itemize}
\raggedright
  \item \emph{Say it:} \emph{mutig}
  \item \emph{Say it:} \emph{mutig}, once more
  \item \emph{Say it:} say \emph{mutig}
  \item \emph{Recall:} say the German for relaxed, then the German for thin, then say \emph{mutig} again
\end{itemize}

\begin{tcolorbox}[breakable,colback=teal!4,colframe=teal!35!black,title={Before you move on}]
What does mutig mean? (\textquotedblleft{}Brave\textquotedblright{}.) Say it once more.
\end{tcolorbox}
